\section{Experimental Setup}
\label{sec:setup}

This section describes the model configuration, anomaly score, operating-point selection, baselines, and evaluation metrics. The main experiments are conducted on CICIDS2017 under the strict no-leakage protocol defined in Section~\ref{sec:protocol}. Supplementary results on SWaT and UNSW-NB15 are reported separately and are used only to provide additional validation context and stress-test evidence.

\subsection{Model Configuration}

The proposed detector is the TCN-WGAN-GP model defined in Section~\ref{sec:method}, evaluated in window-level mode under the strict protocol in Section~\ref{sec:protocol}. Training uses model-train benign windows only, and inference produces one anomaly score per window that is thresholded after independent benign calibration.

For the main CICIDS2017 evaluation, we use one fixed trained model instance selected from the strict training pipeline before final test reporting. To keep the manuscript focused on protocol-level reproducibility rather than implementation logs, we omit filesystem paths in the main text and provide exact checkpoint identifiers and artifact pointers in the released package.

\subsection{TAD Score}

We use the fused TAD score from Section~\ref{sec:method}, combining critic evidence and benign-reference feature deviation. In the selected CICIDS2017 configuration, the fixed mixing coefficient is $\alpha=0.24$, and no post hoc re-tuning is performed on test labels.

\subsection{Operating Point}

The selected CICIDS2017 operating point uses calibration target $\mathrm{FPR}=0.25$. Here, target FPR is a calibration-side quantile parameter on independent benign calibration scores; it is not the claimed deployed benign alarm rate. The threshold selected from the independent benign calibration split is:

\begin{quote}
threshold = 0.425420.
\end{quote}

This threshold is then applied unchanged to the CICIDS2017 test split. The observed benign false-positive rate on the test split is 0.0460. Throughout the paper, we distinguish the calibration target FPR from the observed test benign FPR. The former defines the threshold-selection rule on development benign data; the latter is the deployment-facing quantity that measures realized alert noise on unseen traffic. The selected target\_fpr \(=0.25\) is reported because, under this strict calibration rule, it yields high recall while keeping observed test benign FPR below 5\%.

\subsection{Representative Strict-Protocol Baselines}

We compare against representative baselines evaluated under the same window-level strict protocol. These are protocol-compatible baselines used to contextualize the proposed method under a controlled setting, not a comprehensive state-of-the-art comparison.

The evaluated baselines include:

\begin{itemize}
    \item \textbf{IsolationForest}, a classical tree-based anomaly detector;
    \item \textbf{OneClassSVM}, a kernel-based one-class anomaly detector;
    \item \textbf{MLP}, a supervised window-level baseline trained only with training-split window labels;
    \item \textbf{GANomaly}, a representative adversarial/reconstruction-based anomaly detector;
    \item \textbf{TranAD}, a transformer-based time-series anomaly detector.
\end{itemize}

\begin{table}[t]
  \centering
  \caption{Representative baseline grouping used in this study.}
  \label{tab:baseline_grouping}
  \small
  \resizebox{\columnwidth}{!}{%
  \begin{tabular}{llll}
    \toprule
    Method & Family & Training Signal & Role \\
    \midrule
    IsolationForest & Tree one-class & Benign windows & Classical unsupervised baseline \\
    OneClassSVM & Kernel one-class & Benign windows & Classical unsupervised baseline \\
    MLP (Window) & Supervised MLP & Train-split window labels & Supervised reference baseline \\
    GANomaly & Adversarial reconstruction & Benign windows & Deep anomaly baseline \\
    TranAD & Transformer anomaly model & Benign windows & Deep sequence baseline \\
    Ours & TCN-WGAN-GP + TAD & Benign windows + independent calibration & Calibrated auditable detector \\
    \bottomrule
  \end{tabular}%
  }
\end{table}


All baselines are reported using the same window construction and window-level metric definitions. Baselines were selected because they can be aligned to the same strict protocol constraints: benign-only training for anomaly detectors or explicit supervised training assumptions, independent benign calibration, and fixed-threshold evaluation on an untouched test split. We emphasize threshold-dependent metrics in addition to ranking metrics because operational intrusion detection requires a concrete alarm threshold. We also distinguish model families by supervision assumptions: unsupervised methods are trained on model-train benign windows only, while the MLP baseline uses training-split window labels only. This separation keeps comparisons protocol-consistent and avoids hidden label leakage.

\subsection{Metrics}

We report AUC and average precision (AP) to measure ranking quality. We also report observed benign false-positive rate, precision, recall, and F1 score at the selected threshold. Since false positives directly affect analyst workload and alert fatigue, observed benign FPR is treated as a primary deployment-facing metric.

For attack-type analysis, we report Recall@\(\tau\) and mean fused TAD score for each attack family present in the CICIDS2017 test split. This analysis is important because aggregate recall can hide failures on smaller attack categories.

For explanation analysis, we use Integrated Gradients and masking-based decision checks. Attribution-guided masking is compared with random masking to test whether localized evidence regions materially affect alarm scores. This is a post-hoc analyst-facing audit trace, not a complete proof of explanation faithfulness.

\subsection{Supplementary Cross-Dataset Evaluation}

Supplementary cross-dataset evaluation is reported only as scope extension beyond the CICIDS2017 main evidence chain. SWaT is used as supplementary industrial-control validation, and UNSW-NB15 is used as a distribution-shift stress test. These results are interpreted as bounded transfer checks and failure-mode evidence rather than as comprehensive cross-dataset generalization claims.
